\chapter{General Introduction}

\vspace*{1cm}

Artificial Intelligence has become an important technology in many modern
applications, especially with the development of language models capable of
understanding instructions and generating natural language responses. However,
large language models usually require significant computational resources, memory,
and energy. These requirements make their deployment difficult on edge devices,
where hardware capacity is limited.

\vspace*{0.5cm}

Edge AI refers to the execution of artificial intelligence models directly on
local devices such as smartphones, embedded systems, sensors, or small computing
boards. This approach reduces dependence on cloud servers, improves response time,
and can increase privacy because data does not always need to be sent to remote
infrastructure. Nevertheless, deploying AI models on edge devices requires models
that are compact, efficient, and fast during inference.

\vspace*{0.5cm}

For this reason, Small Language Models represent an interesting solution. They are
lighter than large-scale models while still being able to perform useful language
tasks after fine-tuning. In this project, the model \texttt{google/flan-t5-small}
is used as a compact instruction-following model. The objective is to adapt it to
an instruction-tuning dataset and analyze how different optimization algorithms
affect the training process and the final performance.

\vspace*{0.5cm}

The project focuses on the comparison of several optimization methods, including
Stochastic Gradient Descent, SGD with momentum, Adam, and L-BFGS. Newton-CG is
also considered from a theoretical point of view, but it is not implemented because
it is not practical for standard transformer fine-tuning using PyTorch. The
comparison is based not only on model quality, but also on criteria related to
Edge AI, such as training speed, memory consumption, inference latency, model size,
and energy consumption when available.

\vspace*{0.5cm}

The main goal of this report is to present the design and implementation of the
fine-tuning pipeline, explain the optimization algorithms studied, and evaluate
their relevance for efficient Small Language Model adaptation. This work also
serves as a first step toward future compression techniques, such as quantization
or pruning, which can further improve the suitability of language models for Edge
AI applications.
